\section{转移矩阵}
\label{chap:phys-sm-transfer-matrix-fermionization}

图展开把二维 Ising 的配分函数组织成边集之和；要得到完整自由能和有限尺寸谱，更有效的做法是把一个方向的统计求和改写为线性算子的幂。先建立转移算符的一般有限条带结构：精确矩阵元、Perron--Frobenius 主本征值、谱表示的相关函数以及 Pauli 算符分解。随后再引入费米子化，从而把“转移算子存在”与“方格 Ising 恰好能被 Clifford 代数线性化”分开。
\subsection{逐行配置空间与精确转移算子}
\label{sec:phys-sm-transfer-matrix-foundations}

考虑 \(M\times N\) 各向异性方格，水平方向耦合 \(K_h=\beta J_h\)，垂直方向耦合 \(K_v=\beta J_v\)。一行配置记为
\[
 \boldsymbol\sigma=(\sigma_1,\ldots,\sigma_N),
 \qquad \sigma_j=\pm1.
\]
所有行配置张成 \(2^N\) 维配置向量空间
\begin{equation}
 \mathscr R_N=(\mathbb C^2)^{\otimes N}.
 \label{eq:phys-sm-row-space}
\end{equation}
这里的右矢只是经典行配置的线性编码，不是额外的量子时间演化 Hilbert 空间。

横向取周期边界 \(\sigma_{N+1}=\sigma_1\)。把一行内部的水平 Boltzmann 权重平均分到相邻两次纵向传播，定义对称转移算符 \(\mathsf T_N\) 的矩阵元
\begin{equation}
 \langle\boldsymbol\sigma'|\mathsf T_N|\boldsymbol\sigma\rangle
 =\exp\!\left[
 K_v\sum_{j=1}^{N}\sigma_j'\sigma_j
 +\frac{K_h}{2}\sum_{j=1}^{N}
 \left(\sigma_j\sigma_{j+1}
 +\sigma_j'\sigma_{j+1}'\right)
 \right].
 \label{eq:phys-sm-row-transfer-matrix}
\end{equation}
纵向也取周期边界时，有限环面配分函数严格等于
\begin{equation}
 Z_{N,M}=\Tr_{\mathscr R_N}\mathsf T_N^{\,M}.
 \label{eq:phys-sm-transfer-trace-general}
\end{equation}
这个恒等式在任意有限 \(M,N\) 成立，尚未取任何热力学极限。

\subsection{Perron--Frobenius 与长柱面极限}

有限实耦合下，\(\mathsf T_N\) 的每个矩阵元都严格为正。Perron--Frobenius 定理因此给出唯一的最大正本征值 \(\Lambda_0(N)>0\)，且
\[
 |\Lambda_a(N)|<\Lambda_0(N),\qquad a\ge1.
\]
固定宽度 \(N\) 后，
\begin{equation}
 \lim_{M\to\infty}\frac1M\log Z_{N,M}
 =\log\Lambda_0(N),
 \label{eq:phys-sm-transfer-long-cylinder}
\end{equation}
因此每格点有限宽自由能为
\begin{equation}
 -\beta f_N=\frac1N\log\Lambda_0(N).
 \label{eq:phys-sm-strip-free-energy}
\end{equation}
真正的二维自由能还需再取 \(N\to\infty\)。Perron--Frobenius 只控制固定 \(N\) 的有限维谱，并不能单独证明第二个极限。

\begin{proof}
固定行宽 \(N\)。转移矩阵的每个元素都是正的 Boltzmann 权，故有限维 Perron--Frobenius 定理给出单一正本征值 \(\Lambda_0\)，且其余本征值满足 \(|\Lambda_a|<\Lambda_0\)。矩阵在这里还是实对称的，所以可取一组正交本征向量。由迹公式，
\[
Z_{N,M}
=\Tr\mathsf T_N^M
=\Lambda_0^M\left[
1+\sum_{a\ge1}\left(\frac{\Lambda_a}{\Lambda_0}\right)^M
\right].
\]
因求和项数对固定 \(N\) 有限，括号随 \(M\to\infty\) 趋于 \(1\)，从而
\(M^{-1}\log Z_{N,M}\to\log\Lambda_0\)，再除以 \(N\) 得式\cref{eq:phys-sm-strip-free-energy}。此证明没有对 \(N\to\infty\) 给出一致估计，故不能把两个极限偷偷合并。
\end{proof}

最小的算例是一维无场链：每行只留一个自旋、行内不加耦合，令相邻行的耦合为 $K=\beta J>0$。此时
\[
T=\begin{pmatrix}\ee^K&\ee^{-K}\\\ee^{-K}&\ee^K\end{pmatrix},
\qquad
\lambda_+=2\cosh K,\quad\lambda_-=2\sinh K.
\]
直接以向量 $(1,1)^{\mathsf T}$、$(1,-1)^{\mathsf T}$ 对角化，便有 $Z_M=\lambda_+^M+\lambda_-^M$。固定 $K$ 后取 $M\to\infty$，每点的对数配分函数为 $\log(2\cosh K)$；而自旋算符把两个本征向量互换，所以连通二点函数正比于 $(\lambda_-/\lambda_+)^r=(\tanh K)^r$，相关长度为 $1/\log\coth K$。有限宽度的主本征值计算是严格的；再令宽度趋于无穷并寻找二维临界点是另一个极限，不能靠把有限维 Perron--Frobenius 定理换个名字就自动完成。

\subsection{谱表示与依赖观测量的相关长度}

转移谱不只决定自由能，也决定沿纵向的相关衰减。令 \(\widehat O\) 是由一行局域观测量定义的对角算符，取 \(M\to\infty\) 后，间隔 \(r\) 行的连通关联可写成
\begin{equation}
 \langle O_0O_r\rangle_c
 =\sum_{a\ge1}
 C_a(O)
 \left(\frac{\Lambda_a}{\Lambda_0}\right)^r,
 \label{eq:phys-sm-transfer-correlation-spectrum}
\end{equation}
其中
\[
 C_a(O)=
 \langle0|\widehat O|a\rangle
 \langle a|\widehat O|0\rangle
\]
在 \(\mathsf T_N\) 对称实矩阵时非负。于是相关长度不是“第二大本征值”这一句话就完全决定，而要先看观测算符是否与该谱态有非零矩阵元。定义与 \(O\) 耦合的最低非零谱隙
\begin{equation}
 \xi_O^{-1}(N)
 =\min_{a:C_a(O)\neq0}
 \log\frac{\Lambda_0}{|\Lambda_a|}.
 \label{eq:phys-sm-observable-correlation-length}
\end{equation}
这条选择定则在后面区分单费米转移质量与具体自旋相关长度时至关重要。

\begin{proof}
固定 \(N,r\)，在长度为 \(M\) 的周期柱上，精确的有限体积公式是
\[
\langle O_0O_r\rangle_{N,M}
=\frac{\Tr(\widehat O\,\mathsf T_N^r\widehat O\,
\mathsf T_N^{M-r})}{\Tr\mathsf T_N^M}.
\]
把 \(M\to\infty\) 而保持 \(N,r\) 固定。由主本征值的严格占优，
\(\Lambda_0^{-(M-r)}\mathsf T_N^{M-r}\to|0\rangle\langle0|\)，
于是
\[
\langle O_0O_r\rangle_N
=\Lambda_0^{-r}\langle0|\widehat O\mathsf T_N^r\widehat O|0\rangle
=\sum_a\left(\frac{\Lambda_a}{\Lambda_0}\right)^r
|\langle a|\widehat O|0\rangle|^2.
\]
\(a=0\) 项是 \(\langle O\rangle_N^2\)；将它减去便得式\cref{eq:phys-sm-transfer-correlation-spectrum}。若非主谱中有与 \(O\) 耦合的态，取其中最大的 \(|\Lambda_a|\)，其指数衰减包络的倒长度为式\cref{eq:phys-sm-observable-correlation-length}。本征值若为负，相关函数本身还会随 \(r\) 交替变号，因此“衰减包络”比单说“正指数函数”更准确。
\end{proof}

\subsection{Pauli 表示与对偶耦合}
\label{sec:phys-sm-pauli}

在单点基 \(|+\rangle,|-\rangle\) 上定义 Pauli 算符 \(Z|\sigma\rangle=\sigma|\sigma\rangle\) 和翻转算符 \(X|\sigma\rangle=|-\sigma\rangle\)。行内水平因子写成
\begin{equation}
 \mathsf V_h
 =\exp\!\left(K_h\sum_{j=1}^{N}Z_jZ_{j+1}\right).
 \label{eq:phys-sm-horizontal-transfer-factor}
\end{equation}
单点垂直权重矩阵
\[
 \mathsf W_v=
 \begin{pmatrix}
 \ee^{K_v}&\ee^{-K_v}\\
 \ee^{-K_v}&\ee^{K_v}
 \end{pmatrix}
\]
可以写成 \(X\) 的指数。定义对偶耦合 \(K_v^*>0\) 使
\begin{equation}
 \tanh K_v^*=\ee^{-2K_v},
 \qquad
 \sinh(2K_v)\sinh(2K_v^*)=1,
 \label{eq:phys-sm-vertical-dual-coupling}
\end{equation}
则
\begin{equation}
 \mathsf W_v
 =\sqrt{2\sinh(2K_v)}\,\ee^{K_v^*X}.
 \label{eq:phys-sm-vertical-pauli-factor}
\end{equation}
因此
\begin{equation}
 \mathsf T_N
 =\left[2\sinh(2K_v)\right]^{N/2}
 \mathsf V_h^{1/2}\mathsf V_v\mathsf V_h^{1/2},
 \qquad
 \mathsf V_v=\exp\!\left(K_v^*\sum_{j=1}^{N}X_j\right).
 \label{eq:phys-sm-transfer-pauli-factorization}
\end{equation}

\begin{proof}
因 \(X^2=I\)，指数级数的奇、偶项分别相加，得到
\(\ee^{K_v^*X}=\cosh K_v^*\,I+\sinh K_v^*\,X\)。
令 \(\mathsf W_v=A_v\ee^{K_v^*X}\)，比较对角元与非对角元，恰得到
\[
A_v\cosh K_v^*=\ee^{K_v},\qquad
A_v\sinh K_v^*=\ee^{-K_v}.
\]
相除给出 \(\tanh K_v^*=\ee^{-2K_v}\)；平方相减给出
\(A_v^2=\ee^{2K_v}-\ee^{-2K_v}=2\sinh(2K_v)\)。
再用 \(\sinh(2t)=2\tanh t/(1-\tanh^2t)\)，
\[
\sinh(2K_v^*)
=\frac{2\ee^{-2K_v}}{1-\ee^{-4K_v}}
=\frac1{\sinh(2K_v)}.
\]
这证明了单点因子及对偶耦合的两式。各站点的 \(X_j\) 彼此对易，故单点因子的张量积是
\(A_v^N\exp(K_v^*\sum_jX_j)\)。最后把对角的水平权重均分在它两侧，便得到对称转移矩阵式\cref{eq:phys-sm-transfer-pauli-factorization}。
\end{proof}

\eqnrefs{eq:phys-sm-transfer-pauli-factorization} 仍是 \(2^N\) 维谱问题。平移对称只能按总晶格动量分块，并不能把指数维度降到线性规模。真正的可解结构是：沿 $X$ 方向定义 Jordan--Wigner 弦以后，\(\sum X_j\) 与 \(\sum Z_jZ_{j+1}\) 都成为 Majorana 双线性式，而二次 Clifford 算符的共轭作用只在线性 Majorana 空间中演化。因此必须先处理有限环边界项：它决定费米子奇偶与 NS/R 动量网格，不能在取热力学极限之前省略。

先把弦放在转移矩阵的一次 Pauli 方向上，避免将局域翻转项误写成带弦的一次费米子项。以下假定环上至少有三个自旋，定义
\[
S_j=\prod_{k<j}X_k,\qquad
a_j=S_jZ_j,\qquad b_j=S_jY_j,\qquad
P=\prod_{j=1}^N X_j .
\]
空积规定为 \(S_1=I\)。每个 \(a_j,b_j\) 都是厄米算符且平方为 \(I\)。同一位置上 \(Z_jY_j=-Y_jZ_j\)；若 \(j<k\)，较长的弦 \(S_k\) 恰包含一个 \(X_j\)，它与位置 \(j\) 的 \(Z_j\) 或 \(Y_j\) 反对易，而其余因子对易。因此
\begin{equation}
\{a_j,a_k\}=2\delta_{jk},\qquad
\{b_j,b_k\}=2\delta_{jk},\qquad
\{a_j,b_k\}=0 .
\label{eq:phys-sm-majorana-relations}
\end{equation}
这才是 \(2N\) 个 Majorana 生成元的 Clifford 关系。取与下一节升降算符一致的约定
\(c_j=(a_j-\ii b_j)/2\)。这些关系随即给出
\(\{c_j,c_k^\dagger\}=\delta_{jk}\) 与 \(\{c_j,c_k\}=0\)。同站点计算
\(c_j^\dagger c_j=(1-\ii a_jb_j)/2=(1-X_j)/2\)，所以
\(P=\prod_jX_j=(-1)^{\sum_jc_j^\dagger c_j}\)。

关键是原转移因子中的两个局域算符都恰为 Majorana 的\emph{双线性}。使用 \(Z_jY_j=-\ii X_j\) 和 \(Y_jX_j=-\ii Z_j\)，逐项相乘可得
\begin{equation}
X_j=\ii a_jb_j,\qquad
Z_jZ_{j+1}=\ii b_ja_{j+1}\quad(1\le j<N).
\label{eq:phys-sm-majorana-bulk}
\end{equation}
环的最后一条边不能照抄第二式。因为 \(S_N=P X_N\)，有
\(b_Na_1=S_NY_NZ_1=\ii PZ_NZ_1\)，故
\begin{equation}
Z_NZ_1=-P\,\ii b_Na_1.
\label{eq:phys-sm-majorana-boundary}
\end{equation}
于是前面的 \(\mathsf V_v=\exp(K_v^*\sum_jX_j)\) 和
\(\mathsf V_h=\exp(K_h\sum_jZ_jZ_{j+1})\)
在固定 \(P\) 的子空间中都成为二次 Majorana 算符的指数。

这里必须分清“单个费米子”与“偶次费米子式”。\(P\) 与每个 \(a_j,b_j\) \emph{反}对易，故与 \(c_j,c_j^\dagger\) 也反对易；它只与式\eqnrefs{eq:phys-sm-majorana-bulk,eq:phys-sm-majorana-boundary}中的双线性算符对易。因此转移矩阵保留两个宇称子空间，却不能在整个 Fock 空间上把 \(P\) 当作无条件的数。固定 \(P=+1\) 时，最后一条键相对于体内键多负号，对应反周期 Majorana 边界；固定 \(P=-1\) 时则对应周期边界。若改变 Majorana 或宇称的定义，NS/R 名称也可能对调，但式\cref{eq:phys-sm-majorana-boundary}的算符符号必须先算清楚。

有限环的转移矩阵仍是两个不对易指数的乘积，必须先在每个宇称扇区作 Fourier 变换，并单独处理动量 \(0,\pi\) 的特殊模式。这样才能得到正确的单粒子快度
\begin{equation}
\cosh\epsilon(k)=
\cosh(2K_h)\cosh(2K_v^*)
-\sinh(2K_h)\sinh(2K_v^*)\cos k.
\label{eq:phys-sm-fermion-dispersion}
\end{equation}
它在 \(k=0\) 的质量是 \(2|K_h-K_v^*|\)，仅当 \(K_h=K_v^*\) 闭合。这里的快度属于两个\emph{不对易指数}的乘积，必须先在 Majorana 空间中复合两次共轭作用，再求本征值。若先把两个指数的生成元相加后对角化，所得是另一个算符的谱，通常不等于此处的转移快度。

转移矩阵与一维量子链的关系要在\emph{各向异性连续极限}中建立，而不能把任意固定耦合下两个不对易指数的乘积直接写成指数之和。将式\cref{eq:phys-sm-transfer-pauli-factorization}中的整体正标量记为 $A_N$，并取
\[
K_h=\varepsilon J,\qquad K_v^*=\varepsilon h,
\qquad \varepsilon\downarrow0,
\]
其中 $J,h>0$ 固定。由 $\tanh K_v^*=\ee^{-2K_v}$，这同时意味着原二维模型的竖直耦合 $K_v\to\infty$，故确实是强各向异性而非普通热力学极限。归一化后的对称转移步为
\[
A_N^{-1}\mathsf T_N
=\exp\!\left(\frac{\varepsilon J}{2}\sum_jZ_jZ_{j+1}\right)
 \exp\!\left(\varepsilon h\sum_jX_j\right)
 \exp\!\left(\frac{\varepsilon J}{2}\sum_jZ_jZ_{j+1}\right).
\]
固定环长 $N$ 时这是有限维矩阵，逐项展开便得
\[
A_N^{-1}\mathsf T_N=I-\varepsilon H_{\rm Q}+O(\varepsilon^2),
\qquad
H_{\rm Q}=-J\sum_jZ_jZ_{j+1}-h\sum_jX_j.
\]
若令转移步数 $M\to\infty$、$M\varepsilon\to t$，矩阵乘积极限或 Trotter 公式给出
$\lim(A_N^{-1}\mathsf T_N)^M=\ee^{-tH_{\rm Q}}$。这时、也只有在这项受控极限中，二维经典模型的转移方向才变成一维横场 Ising 链的虚时间。二维自对偶条件 $K_h=K_v^*$ 在这个极限中化为 $J=h$。先固定 $N$ 建立矩阵极限，再研究 $N\to\infty$ 的量子临界行为；两个极限不能在未经估计时交换。
